\documentclass[11pt,a4paper]{article}

\usepackage[utf8]{inputenc}
\usepackage[T1]{fontenc}
\usepackage[french]{babel}
\usepackage{lmodern}
\usepackage{geometry}
\usepackage{amsmath,amssymb}
\usepackage{booktabs}
\usepackage{array}
\usepackage{hyperref}
\usepackage{xcolor}

\geometry{margin=2.5cm}
\hypersetup{
  colorlinks=true,
  linkcolor=blue!45!black,
  urlcolor=blue!45!black,
  citecolor=blue!45!black
}

\newcommand{\STFT}{\operatorname{STFT}}
\newcommand{\argmin}{\operatorname*{arg\,min}}
\newcommand{\argmax}{\operatorname*{arg\,max}}
\newcommand{\code}[1]{\texttt{\detokenize{#1}}}

\title{Rapport sur un algorithme de separation de sources acoustiques}
\author{A completer}
\date{\today}

\begin{document}

\maketitle

\begin{abstract}
Ce rapport decrit un pipeline de traitement audio multicanal destine a isoler
plusieurs sources acoustiques puis a estimer leurs differences de temps
d'arrivee entre microphones. Le pipeline s'appuie d'abord sur une selection de
bins dans le plan temps-frequence, afin de ne conserver que les zones contenant
du signal exploitable. La partie principale du rapport porte ensuite sur
l'algorithme de separation de sources implemente dans le module
\code{BSS}. Cette methode modelise chaque bin actif comme un coefficient
complexe de source multiplie par un vecteur spatial de retards, applique un
clustering EM independant par frequence, puis resout le probleme de permutation
entre frequences par un RANSAC circulaire sur les phases relatives des
centroïdes. Enfin, les sources separees ou les bins isoles peuvent etre
utilises pour estimer les TDOA, notamment via des methodes de beamforming
masquees ou ponderees.
\end{abstract}

\tableofcontents
\newpage

\section{Introduction}

L'estimation de la position ou de la direction d'une source acoustique dans un
milieu naturel repose souvent sur l'exploitation d'un reseau de microphones. Si
une source emet un signal, celui-ci est observe par chaque microphone avec un
retard et une attenuation qui dependent de la geometrie de propagation. Les
differences de temps d'arrivee, ou TDOA pour \emph{Time Difference Of Arrival},
contiennent donc une information spatiale directement liee a la source.

Dans un contexte realiste, plusieurs sources peuvent etre actives
simultanement, et le signal observe est un melange de contributions utiles et de
bruits. Le probleme consiste alors a isoler les composantes associees a chaque
source avant d'estimer leurs retards respectifs. Le pipeline considere dans ce
rapport est organise en trois etapes principales :

\begin{enumerate}
  \item construire un masque temps-frequence afin de selectionner les bins ou
  du signal est present ;
  \item separer les sources a l'interieur de ces bins actifs ;
  \item calculer les TDOA a partir des bins isoles, en particulier par
  beamforming.
\end{enumerate}

La contribution principale de ce rapport concerne la deuxieme etape. Elle
correspond au module de separation de sources present dans \code{BSS}, et plus
precisement a l'algorithme de type Sawada/LOST implemente dans
\code{BSS/Algo_Separation/Sawada_separation.py}. Cette methode exploite la
structure spatiale des coefficients de STFT multicanaux : dans une zone
temps-frequence locale, un bin actif est souvent domine par une seule source.
Les observations associees a une meme source forment alors des directions
complexes coherentes dans l'espace des microphones.

\section{Creation d'un masque temps-frequence }
\subsection{Representation temps-frequence multicanale}

On note $x_m[n]$ le signal temporel observe par le microphone $m$, avec
$m \in \{1,\dots,M\}$. Pour chaque microphone, une transformee de Fourier a
court terme est calculee :
\[
  X_m(f,t) = \STFT\{x_m[n]\}(f,t),
  \]
  ou $f$ designe l'indice de frequence et $t$ l'indice de trame temporelle.
  
  Les observations multicanales sont ensuite regroupees sous forme vectorielle :
  \[
    \mathbf{x}_{f,t}
    =
    \begin{bmatrix}
      X_1(f,t) & X_2(f,t) & \cdots & X_M(f,t)
    \end{bmatrix}^{\top}
    \in \mathbb{C}^{M}.
\]

Cette representation est centrale pour le pipeline. L'amplitude des
coefficients permet d'identifier les zones temps-frequence contenant de
l'energie utile. La phase relative entre les microphones permet quant a elle
d'exploiter la geometrie de propagation et donc de distinguer les sources.



Le masque temps-frequence a pour objectif de selectionner les bins $(f,t)$ qui
contiennent du signal utile. Cette etape reduit le nombre d'observations
traitees par l'algorithme de separation et limite l'influence des bins domines
par le bruit.



\subsection{Selection des bins actifs par energie}

 la STFT multicanale brute est d'abord
calculee. L'energie d'un bin temps-frequence est ensuite definie par :
\[
  E(f,t) = \sum_{m=1}^{M} |X_m(f,t)|^2.
\]

Cette energie est convertie en decibels :
\[
  E_{\mathrm{dB}}(f,t)
    =
    10 \log_{10}\left(E(f,t)+\epsilon\right).
\]

Le seuil d'activation est calcule relativement a un plancher d'energie estime
par percentile :
\[
  \theta
    =
    \operatorname{percentile}_{p}\left(E_{\mathrm{dB}}\right)
    +
    \Delta_{\mathrm{dB}}.
\]

Un bin est conserve si :
\[
  E_{\mathrm{dB}}(f,t) \geq \theta.
\]

Dans la configuration de benchmark presente dans
\code{BSS/Benchmark/runners.py}, les valeurs utilisees sont typiquement un
plancher au percentile $20$, un seuil place $6$ dB au-dessus de ce plancher, et
un minimum de $3$ trames actives consecutives par frequence.

Ensuite,
un bin actif isole ne fournit pas assez d'information locale pour stabiliser le
clustering EM. L'implementation filtre donc les activations frequence par
frequence en conservant seulement les groupes temporels suffisamment longs.
Si $L_{\min}$ designe la longueur minimale, un segment actif n'est conserve que
si :
\[
  \#\{t : (f,t) \text{ actif dans le segment}\} \geq L_{\min}.
\]

Ce filtrage est important pour la robustesse de la separation. En pratique, il
evite que des fluctuations locales du bruit soient presentees a l'EM comme des
observations fiables.

On pourrait également penser à utiliser des filtres plus devellopés voir meme des CNN. Cela nous mène à la 
seconde approche. 

\subsection{Utilisation d'un réseau de neuronnes}
\section{Separation de sources dans le masque temps-frequence}

\subsection{Préliminaires}
\subsubsection{Hypothese de parcimonie locale}

L'algorithme de separation repose sur une hypothese de parcimonie
temps-frequence. Dans une zone locale du plan temps-frequence, un bin actif est
souvent domine par une seule source. Pour une source $k$, on peut alors ecrire :
\[
  \mathbf{x}_{f,t}
    \simeq
    s_k(f,t)\,\mathbf{a}_k(f),
\]
ou $s_k(f,t) \in \mathbb{C}$ est le coefficient complexe de la source et
$\mathbf{a}_k(f) \in \mathbb{C}^M $ est le vecteur spatial associe a cette source.

Le probleme de separation devient donc un probleme de regroupement de
directions complexes. A frequence fixee, les bins domines par la meme source
doivent se regrouper autour d'un meme vecteur spatial.

\subsubsection{Justification de la forme des bins}

On suppose que la propagation d'une source $k$ vers un microphone $m$ peut etre
approximativement decrite par un gain $\rho_{m,k}$ et un retard
$\tau_{m,k}$. Dans le domaine frequentiel, un retard temporel se traduit par
une phase lineaire en frequence :
\[
  X_m(f,t)
    \simeq
    \rho_{m,k}
    e^{-j2\pi f\tau_{m,k}}
    S_k(f,t).
\]

En empilant les observations de tous les microphones, on obtient :
\[
  \mathbf{x}_{f,t}
    \simeq
    S_k(f,t)
    \begin{bmatrix}
      \rho_{1,k} e^{-j2\pi f\tau_{1,k}} \\
      \rho_{2,k} e^{-j2\pi f\tau_{2,k}} \\
      \vdots \\
      \rho_{M,k} e^{-j2\pi f\tau_{M,k}}
    \end{bmatrix}
    =
    S_k(f,t)\,\mathbf{a}_k(f).
\]

Cette equation justifie la forme utilisee par l'algorithme : un bin est modelise
comme un coefficient complexe scalaire, propre au contenu de la source, multiplie
par un vecteur complexe qui encode la signature spatiale. Cette signature
contient en particulier les retards relatifs entre microphones. De plus, comme les micros sont proches,
on peut neglecter les differences d'amplitude entre les microphones.
 (et surtout ne pas s'appuyer dessus pour l'algorithme)

\subsubsection{Normalisation et blanchiment des observations}

L'algorithme cherche principalement a separer les directions spatiales, et non
les amplitudes. Les bins sont donc normalises individuellement :
\[
  \tilde{\mathbf{x}}_{f,t}
    =
    \frac{\mathbf{x}_{f,t}}
         {\|\mathbf{x}_{f,t}\|_2 + \epsilon}.
\]

Apres normalisation, on obtient approximativement :
\[
  \tilde{\mathbf{x}}_{f,t}
    \approx
    e^{j\phi_{f,t}}
    \frac{\mathbf{a}_k(f)}{\|\mathbf{a}_k(f)\|_2}.
\]

Le facteur $e^{j\phi_{f,t}}$ correspond a une phase globale. Il ne porte pas
l'information spatiale utile, qui est contenue dans les phases relatives entre
microphones. C'est pourquoi l'algorithme travaille sur des distances
directionnelles invariantes a cette phase globale.

Dans l'implementation, il est aussi possible d'ajouter une etape de blanchiment
spatial apres cette premiere normalisation. Pour chaque frequence $f$, on
considere la matrice des observations normalisees :
\[
  \tilde{\mathbf{X}}_f
    =
    \begin{bmatrix}
      \tilde{\mathbf{x}}_{f,1} &
      \tilde{\mathbf{x}}_{f,2} &
      \cdots &
      \tilde{\mathbf{x}}_{f,T}
    \end{bmatrix}.
\]

On estime alors la matrice de correlation spatiale :
\[
  \mathbf{C}_f
    =
    \frac{1}{T}
    \tilde{\mathbf{X}}_f
    \tilde{\mathbf{X}}_f^{H}.
\]

Cette matrice est diagonalisee :
\[
  \mathbf{C}_f
    =
    \mathbf{U}_f
    \mathbf{\Lambda}_f
    \mathbf{U}_f^{H},
\]
ou $\mathbf{\Lambda}_f$ contient les valeurs propres et $\mathbf{U}_f$ les
vecteurs propres. La matrice de blanchiment appliquee dans le code est :
\[
  \mathbf{W}_f
    =
    \mathbf{\Lambda}_f^{-1/2}
    \mathbf{U}_f^{H}.
\]

Les observations blanchies sont donc :
\[
  \mathbf{y}_{f,t}
    =
    \mathbf{W}_f
    \tilde{\mathbf{x}}_{f,t}.
\]

Apres blanchiment, les directions sont transformees dans un espace ou la
correlation spatiale globale est normalisee :
\[
  \frac{1}{T}
  \sum_t
  \mathbf{y}_{f,t}\mathbf{y}_{f,t}^{H}
  \approx
  \mathbf{I}.
\]

Cette operation reduit les correlations communes entre microphones et peut
rendre les directions des sources plus faciles a separer par l'EM. En revanche,
les centroïdes estimes ne vivent plus directement dans l'espace physique
original des microphones : ils sont exprimes dans l'espace blanchi. Lorsque
l'on veut interpreter leurs phases comme des retards, il faut donc appliquer
l'inverse du blanchiment aux centroïdes avant l'assignation spatiale ou
l'analyse des TDOA. C'est ce que fait la propriete
\code{all_centroids_unwhitened} : si le blanchiment est active, elle applique
la pseudo-inverse de $\mathbf{W}_f$ aux centroïdes avant de les transmettre au
RANSAC circulaire.

Dans \code{SawadaBSS}, les observations sont renormalisees apres le blanchiment.
Ainsi, le clustering EM travaille toujours sur des bins de norme unitaire, que
le blanchiment soit active ou non.


\subsection{Algorithme de SAWADA}
\subsubsection{Clustering EM par frequence}

Pour chaque frequence $f$, l'algorithme extrait les observations actives, on utilise le masque donnée par le réseau de neuronnes ou bien les bins les plus actifs:
\[
  \tilde{\mathbf{x}}_{f,t}, \quad t \in \mathcal{T}_f.
  \]
  
  L'objectif est d'ajuster $K$ clusters, ou $K$ est le nombre de sources. Chaque
  cluster $k$ est represente par trois parametres :
  
  \begin{itemize}
    \item un centroïde complexe $\mathbf{a}_{f,k}$ ;
    \item une variance $\sigma^2_{f,k}$, qui mesure la dispersion autour de la
    direction ;
    \item un poids de melange $\alpha_{f,k}$.
  \end{itemize}
  
  L'initialisation est faite par K-means sur les amplitudes des observations, puis
  les centroïdes sont normalises sur la sphere unite complexe.
  
  \subsubsection{Etape E}
  
  L'etape E calcule, pour chaque observation, la probabilite posterieure
  d'appartenir a chaque cluster. La distance directionnelle utilisee est :
  \[
    d^2_{k,t}
    =
    1 -
    \left|
    \mathbf{a}_{f,k}^{H}
    \tilde{\mathbf{x}}_{f,t}
    \right|^2.
\]

Cette distance est faible lorsque l'observation est proche de la direction du
centroïde, independamment d'une phase globale. La probabilite non normalisee
associee au cluster $k$ est :
\[
  q_{k,t}
    =
    \alpha_{f,k}
    \frac{1}
         {\left(\pi(\sigma^2_{f,k}+\epsilon)\right)^{M-1}}
    \exp\left(
      -\frac{d^2_{k,t}}{\sigma^2_{f,k}+\epsilon}
    \right).
\]

Les responsabilites EM sont ensuite obtenues par normalisation :
\[
  \gamma_{k,t}
    =
    \frac{q_{k,t}}
         {\sum_{\ell=1}^{K} q_{\ell,t}+\epsilon}.
\]

\subsubsection{Etape M}

L'etape M met a jour les parametres des clusters a partir des responsabilites.
Les poids de melange sont mis a jour avec un prior de Dirichlet $\phi$ :
\[
  \alpha_{f,k}
    =
    \frac{\sum_t \gamma_{k,t}+\phi}
         {N_f+K\phi},
\]
ou $N_f$ est le nombre d'observations actives a la frequence $f$.

Le centroïde est estime comme la direction principale de la matrice de
correlation ponderee :
\[
  \mathbf{R}_{f,k}
    =
    \sum_t
    \gamma_{k,t}
    \tilde{\mathbf{x}}_{f,t}
    \tilde{\mathbf{x}}_{f,t}^{H}.
\]

Le nouveau centroïde $\mathbf{a}_{f,k}$ est le vecteur propre principal de
$\mathbf{R}_{f,k}$. La variance est ensuite mise a jour comme la dispersion
moyenne ponderee des observations autour de cette direction :
\[
  \sigma_{f,k}^2
    =
    \frac{
      \sum_t \gamma_{k,t}
      \left(
        1-\left|\mathbf{a}_{f,k}^{H}\tilde{\mathbf{x}}_{f,t}\right|^2
      \right)
    }
    {
      (M-1)\sum_t \gamma_{k,t}
    }.
\]

Les demonstrations detaillees de ces trois mises a jour sont donnees en
annexe~\ref{annexe:em-mstep}.

En pratique, ces deux etapes sont alternees un nombre fixe de fois. Dans
l'implementation actuelle. Pour 20 itérations, cela fonctionne bien en pratique.

\subsubsection{Construction des masques de clusters}

Apres convergence de l'EM, un masque binaire est construit par decision dure.
Pour chaque bin actif, on choisit le cluster avec le plus grand score :
\[
  \hat{k}(f,t)
    =
    \argmax_k
    \frac{
      \left|
        \mathbf{a}_{f,k}^{H}
        \tilde{\mathbf{x}}_{f,t}
      \right|^2
    }
    {\sigma^2_{f,k}+\epsilon}.
\]

Le masque du cluster $k$ est alors :
\[
  M_k(f,t)
    =
    \begin{cases}
      1 & \text{si } \hat{k}(f,t)=k,\\
      0 & \text{sinon.}
    \end{cases}
\]

A ce stade, les masques sont coherents localement a une frequence donnee, mais
pas encore globalement entre frequences. Comme l'EM est lance independamment
pour chaque frequence, le cluster numerote $0$ a une frequence peut correspondre
a une autre source que le cluster numerote $0$ a la frequence suivante. C'est le
probleme classique de permutation.

\subsection{Assignation des clusters à l'une des sources}

\subsubsection{Motivation}

Pour reconstruire une source, il faut que les masques soient coherents sur
toutes les frequences. L'implementation propose deux strategies d'alignement :
une methode par correlation d'enveloppes, et une methode par RANSAC circulaire.
La configuration de benchmark utilise l'alignement par RANSAC, car il exploite
directement la structure physique des phases relatives.
\subsubsection{Alignement par correlation d'enveloppes}
L'assignation par correlation d'enveloppes consiste a associer, pour chaque
frequence, les clusters locaux aux sources globales deja suivies aux frequences
precedentes. Pour ce faire, on calcule l'enveloppe $e_k^f(t)$ du cluster $k$ a
la frequence $f$ :
\[
  e_k^f(t)
    =
    \sum_m |X_m(f,t)|^2 M_k(f,t)
\] 
où $M_k(f,t)$ est le masque binaire du cluster $k$ pour la frequence $f$.

On calcule ensuite pour chaque combinaison d'attribution possible 
\[
  \pi \in \mathfrak{S}_K
\]
un score de coherence entre les enveloppes du bin courant et les enveloppes de
reference deja associees aux sources. Si $r_s^{f-1}(t)$ designe l'enveloppe de
reference de la source $s$ construite a partir des frequences precedentes, le
score de la permutation $\pi$ est :
\[
  C_f(\pi)
    =
    \sum_{s=1}^{K}
    \rho\left(
      e_{\pi(s)}^f(t),
      r_s^{f-1}(t)
    \right),
\]
ou $\rho(\cdot,\cdot)$ est la correlation de Pearson :
\[
  \rho(u,v)
    =
    \frac{
      \sum_t
      \left(u(t)-\bar{u}\right)
      \left(v(t)-\bar{v}\right)
    }
    {
      \sqrt{\sum_t \left(u(t)-\bar{u}\right)^2}
      \sqrt{\sum_t \left(v(t)-\bar{v}\right)^2}
      +
      \epsilon
    }.
\]

La permutation retenue est celle qui maximise cette correlation totale :
\[
  \pi_f^{*}
    =
    \argmax_{\pi \in \mathfrak{S}_K}
    C_f(\pi).
\]

Les masques de la frequence $f$ sont ensuite reordonnes selon $\pi_f^{*}$ :
\[
  M_s(f,t)
    \leftarrow
    M_{\pi_f^{*}(s)}(f,t).
\]

L'enveloppe de reference est finalement mise a jour de facon progressive. Dans
le mode \code{ema} utilise par l'implementation, cette mise a jour prend la
forme :
\[
  r_s^{f}(t)
    =
    (1-\beta)r_s^{f-1}(t)
    +
    \beta e_{\pi_f^{*}(s)}^f(t),
\]
avec $\beta=0.1$ dans le code. Une frequence dont l'energie associee a une
source est trop faible ne met pas a jour la reference, afin d'eviter qu'un bin
peu informatif degrade l'alignement des frequences suivantes.

Cette methode est simple et ne suppose pas explicitement un modele de retard.
Elle utilise seulement la similarite temporelle des enveloppes : deux clusters
appartenant a la meme source doivent presenter des activations temporelles
corrélées d'une frequence a l'autre. En revanche, cette méthode ordonne les fréquences car les basses fréquences 
servent initialement à construire l'enveloppe de réference. Aussi, beaucoup de situation, notamment lors de 
vocalise simultané, peuvent donner lieu à des erreurs. 

\subsubsection{RANSAC circulaire} 
\paragraph{conversion des centroides en phase relative} ~\\
Les centroïdes complexes $\mathbf{a}_{f,k}$ contiennent une phase globale
arbitraire. Pour la supprimer, on choisit un microphone de reference, note ici
$1$, et on calcule :
\[
  z^{\mathrm{rel}}_{f,k,m}
    =
    a_{f,k,m}
    a^{*}_{f,k,1}.
\]

La phase relative associee est :
\[
  \theta_{f,k,m}
    =
    \arg\left(z^{\mathrm{rel}}_{f,k,m}\right).
\]

Si le centroïde correspond a une source physique dominee par des retards, alors
on attend une relation de la forme :
\[
  \theta_{f,k,m}
    \approx
    b_m
    -
    2\pi f(\tau_m-\tau_1)
    \pmod{2\pi}.
\]

Les phases relatives d'une meme source decrivent donc une droite en frequence,
mais observee modulo $2\pi$. Cette propriete est la base du RANSAC circulaire
implemente dans \code{BSS/CircularRansac}.

\textit{Remarque : si l'on a précedemment blanchis le données, il faut penser à faire l'opération inverse
sur les centroides au risque de faire un un ransac sur des données qui ne représente aucunement des droites}

\paragraph{Modele ajusté par RANSAC circulaire} ~\\

Le modele ajuste est une droite circulaire multidimensionnelle :
\[
  \boldsymbol{\theta}_i
    =
    \mathbf{b}
    +
    (f_i-f_{\mathrm{ref}})\mathbf{v}
    \pmod{2\pi}.
\]

La distance d'une observation au modele est calculee apres repliement de phase :
\[
  d_i
    =
    \left\|
      \operatorname{wrap}
      \left(
        \boldsymbol{\theta}_i
        -
        \mathbf{b}
        -
        (f_i-f_{\mathrm{ref}})\mathbf{v}
      \right)
    \right\|_2.
\]

Concretement, chaque observation donnee au RANSAC correspond a un centroïde
disponible a une frequence donnee. Elle est decrite par sa frequence $f_i$ et
par son vecteur de phases relatives :
\[
  \boldsymbol{\theta}_i
    =
    \begin{bmatrix}
      \theta_{i,1} & \cdots & \theta_{i,M}
    \end{bmatrix}^{\top}.
\]

Le RANSAC construit d'abord des hypotheses de trajectoires a partir de petits
echantillons d'observations, typiquement deux centroïdes pris a deux frequences
distinctes. Pour une paire $(i,j)$, la pente candidate doit verifier :
\[
  \boldsymbol{\theta}_j-\boldsymbol{\theta}_i
    =
    (f_j-f_i)\mathbf{v}
    \pmod{2\pi}.
\]

Comme les phases sont definies modulo $2\pi$, plusieurs pentes peuvent expliquer
la meme paire d'observations :
\[
  \mathbf{v}
    =
    \frac{
      \boldsymbol{\theta}_j-\boldsymbol{\theta}_i
      +
      2\pi\mathbf{n}
    }
    {f_j-f_i},
    \qquad
    \mathbf{n}\in\mathbb{Z}^{M}.
\]

L'implementation enumere les hypotheses compatibles avec les bornes physiques
sur la pente. Ces bornes proviennent du retard maximal possible entre
microphones : comme $v_m=-2\pi\tau_m$, limiter les TDOA revient a limiter les
pentes admissibles.

Pour chaque hypothese $(\mathbf{b},\mathbf{v})$, l'algorithme calcule la
distance circulaire de toutes les observations au modele. Une observation est
consideree comme un inlier si :
\[
  d_i < \eta,
\]
ou $\eta$ est le seuil de residu circulaire. Le score utilise est de type
MSAC : au lieu de compter seulement les inliers, on somme des residus tronques :
\[
  S
    =
    \sum_i
    \min(d_i^2,\eta^2).
\]

Le meilleur modele est celui qui minimise ce score. Apres selection d'une bonne
hypothese, une optimisation locale raffine la pente et l'intercept sur les
inliers. Pour une pente fixee, l'intercept optimal est obtenu par moyenne
circulaire des phases residuelles :
\[
  b_m(\mathbf{v})
    =
    \arg
    \sum_{i\in\mathcal{I}}
    \exp\left(
      j\left(\theta_{i,m}-(f_i-f_{\mathrm{ref}})v_m\right)
    \right),
\]
ou $\mathcal{I}$ designe l'ensemble des inliers.

Comme il y a plusieurs sources, l'extraction est sequentielle. Le RANSAC trouve
d'abord la trajectoire la plus coherente, associe les centroïdes compatibles a
cette source, puis retire ou marque ces centroïdes avant de chercher la
trajectoire suivante. Cette procedure est repetee jusqu'a obtenir $K$
trajectoires, c'est-a-dire une trajectoire de phase par source attendue.

Enfin, tous les centroïdes disponibles sont assignes a la trajectoire la plus
proche. Les centroïdes trop eloignes de toutes les trajectoires peuvent etre
marques comme non assignes. Les labels obtenus permettent ensuite de fusionner
les masques de clusters en masques de sources coherents sur l'ensemble des
frequences.

\paragraph{Lien entre pente de phase et TDOA} ~\\ 

La pente estimee pour le microphone $m$ a une interpretation directe :
\[
  v_m = -2\pi(\tau_m-\tau_{\mathrm{ref}}).
\]

Le retard relatif au microphone de reference est donc :
\[
  \tau_m-\tau_{\mathrm{ref}}
    =
    -\frac{v_m}{2\pi}.
\]

Les TDOA entre paires de microphones sont ensuite calcules par difference :
\[
  \Delta\tau_{i,j}
    =
    \tau_j-\tau_i.
\]

En revanche, les tdoas estimé à partir de la pente recupérée dans 
le RANSAC sont souvent mauvais en comparaison des autres methodes. Cela est donc utile à améliorer la cohérence 
des mask mais ne sert pas directement au calcul de TDOAS. Il servent principalement à verifier que l'on ne s'est
pas totalement trompé. 

\subsection{Fusion finale des masques}

Une fois les centroïdes associes a des sources, les masques de clusters sont
fusionnes. Pour une source $s$, le masque final est :
\[
  M_s(f,t)
    =
    \sum_{k:\, label(f,k)=s} M_k(f,t).
\]
Dans nos masques, chaque bin actif est associé à un seul cluster.
La fusion des clusters assignés à une même source produit donc directement un masque binaire de source. 
La reconstruction temporelle consiste ensuite à appliquer ce masque à la STFT multicanale originale,
puis à revenir dans le domaine temporel par STFT inverse. A noter que le signal temporel reconstruit 
ne possède pas vraiment une stft similaire à la stft que l'on obtient avec le masque car la stft masqué ne correspond 
à aucun signal temporel réel. C'est simplement le signal temporel ayant la stft la plus proche de la stft masqué.
Ainsi, il n'y à priori aucune raison que les TDOAS calculé soit les bons. On observe dans les faits que cela donne 
des estimations de TDOAS suffisante. 

\[
  \hat{\mathbf{X}}_s(f,t)
    =
    M_s(f,t)\,\mathbf{X}(f,t),
\]
puis a revenir dans le domaine temporel par STFT inverse.

\section{Calcul des TDOA a partir des bins isoles}

\subsection{TDOA sur sources separees}

Une premiere estimation des TDOA peut etre effectué apres
reconstruction des sources par istft. Pour chaque source separee et chaque paire de
microphones, le retard est estime par correlation temporelle. La convention
utilisee est :
\[
  \Delta\tau_{i,j}
    =
    delay(j)-delay(i).
\]

Cette approche donne une estimation directe sur les signaux temporels
reconstruits. Elle depend cependant de la qualite du masque final : une erreur
de separation ou d'assignation peut deteriorer les signaux reconstruits et donc
les TDOA calcules par correlation.

\subsection{Beamforming masque}

\section{Validation sur un dataset synthétique}


\section{Conclusion}

L'algorithme de separation etudie exploite la structure physique des signaux
multicanaux dans le domaine temps-frequence. Sous l'hypothese qu'un bin actif
est localement domine par une source, l'observation multicanale peut etre
modelisee comme un coefficient complexe multiplie par un vecteur spatial de
retards. Cette structure permet de transformer le probleme de separation en un
probleme de clustering de directions complexes.

Le clustering EM estime, pour chaque frequence, des centroïdes correspondant
aux signatures spatiales locales des sources. Le RANSAC circulaire resout
ensuite le probleme de permutation entre frequences en cherchant des
trajectoires de phase compatibles avec des retards physiques. Les masques finaux
permettent de reconstruire les sources et de calculer des TDOA, soit par
correlation sur les signaux reconstruits, soit a partir des pentes de phase, soit
encore en alimentant un beamformer masque.

La validation sur dataset synthetique fournit un cadre reproductible pour
mesurer les erreurs de TDOA et inspecter les artefacts internes de l'algorithme.
Les prochaines etapes naturelles consistent a completer les resultats
quantitatifs, etudier la sensibilite aux niveaux de bruit et comparer
systematiquement l'apport du masque temps-frequence dans la qualite finale de la
separation et des TDOA.

\appendix

\section{Demonstrations des mises a jour EM}
\label{annexe:em-mstep}

Dans l'algorithme de Sawada, l'etape M consiste a maximiser l'esperance de la
log-vraisemblance complete par rapport aux parametres de chaque source. Pour une
frequence donnee, on note
$\mathcal{X}=\{\tilde{\mathbf{x}}_t\}_{t=1}^{T}$ l'ensemble des observations
normalisees, avec $\|\tilde{\mathbf{x}}_t\|_2=1$. Pour une source $i$, les
parametres sont le centroïde $\mathbf{a}_i$, contraint par
$\|\mathbf{a}_i\|_2=1$, la variance $\sigma_i^2$ et le poids de melange
$\alpha_i$.

La quantite maximisee par l'EM est :
\[
  Q
    =
    \sum_{t=1}^{T}
    \sum_{i=1}^{K}
    \gamma_{i,t}
    \left[
      \ln(\alpha_i)
      +
      \ln p\left(
        \tilde{\mathbf{x}}_t
        \mid
        \mathbf{a}_i,\sigma_i^2
      \right)
    \right].
\]

\subsection{Mise a jour des poids}

Sans regularisation, la mise a jour des poids maximise :
\[
  \sum_{i=1}^{K}
  \sum_{t=1}^{T}
  \gamma_{i,t}\ln(\alpha_i)
  \quad
  \text{sous la contrainte}
  \quad
  \sum_{i=1}^{K}\alpha_i=1.
\]

En utilisant un multiplicateur de Lagrange $\lambda$, on derive :
\[
  \frac{\partial}{\partial \alpha_i}
  \left(
    \sum_{i,t}\gamma_{i,t}\ln(\alpha_i)
    -
    \lambda\sum_i\alpha_i
  \right)
  =
  \frac{\sum_t\gamma_{i,t}}{\alpha_i}
  -
  \lambda
  =
  0.
\]

On obtient donc :
\[
  \alpha_i
    =
    \frac{\sum_t\gamma_{i,t}}{\lambda}.
\]

Puisque $\sum_i\alpha_i=1$, on a :
\[
  \lambda
    =
    \sum_i\sum_t\gamma_{i,t}
    =
    T.
\]

D'ou :
\[
  \alpha_i
    =
    \frac{\sum_t\gamma_{i,t}}{T}.
\]

L'ajout de $\phi$ dans le code correspond a l'introduction d'un prior de
Dirichlet pour regulariser l'estimation. Dans ce cas, la mise a jour devient :
\[
  \alpha_i
    =
    \frac{\sum_t\gamma_{i,t}+\phi}
         {T+K\phi}.
\]

\subsection{Mise a jour du centroïde}

Le terme de la log-vraisemblance dependant du centroïde $\mathbf{a}_i$ est :
\[
  Q(\mathbf{a}_i)
    =
    \sum_{t=1}^{T}
    \gamma_{i,t}
    \left(
      -
      \frac{
        \left\|
          \tilde{\mathbf{x}}_t
          -
          \left(\mathbf{a}_i^{H}\tilde{\mathbf{x}}_t\right)\mathbf{a}_i
        \right\|_2^2
      }
      {\sigma_i^2}
    \right).
\]

Puisque $\|\tilde{\mathbf{x}}_t\|_2=1$ et $\|\mathbf{a}_i\|_2=1$, la distance a
la projection verifie :
\[
  \left\|
    \tilde{\mathbf{x}}_t
    -
    \left(\mathbf{a}_i^{H}\tilde{\mathbf{x}}_t\right)\mathbf{a}_i
  \right\|_2^2
  =
  1
  -
  \left|
    \mathbf{a}_i^{H}\tilde{\mathbf{x}}_t
  \right|^2.
\]

Maximiser $Q(\mathbf{a}_i)$ revient donc a maximiser :
\[
  \mathcal{J}(\mathbf{a}_i)
    =
    \sum_{t=1}^{T}
    \gamma_{i,t}
    \left|
      \mathbf{a}_i^{H}\tilde{\mathbf{x}}_t
    \right|^2.
\]

Cette quantite peut s'ecrire sous forme quadratique :
\[
  \mathcal{J}(\mathbf{a}_i)
    =
    \sum_{t=1}^{T}
    \gamma_{i,t}
    \mathbf{a}_i^{H}
    \tilde{\mathbf{x}}_t
    \tilde{\mathbf{x}}_t^{H}
    \mathbf{a}_i
    =
    \mathbf{a}_i^{H}
    \left(
      \sum_{t=1}^{T}
      \gamma_{i,t}
      \tilde{\mathbf{x}}_t
      \tilde{\mathbf{x}}_t^{H}
    \right)
    \mathbf{a}_i.
\]

En posant :
\[
  \mathbf{R}_i
    =
    \sum_{t=1}^{T}
    \gamma_{i,t}
    \tilde{\mathbf{x}}_t
    \tilde{\mathbf{x}}_t^{H},
\]
on cherche :
\[
  \max_{\mathbf{a}_i}
  \mathbf{a}_i^{H}
  \mathbf{R}_i
  \mathbf{a}_i
  \quad
  \text{sous la contrainte}
  \quad
  \mathbf{a}_i^{H}\mathbf{a}_i=1.
\]

Par la methode des multiplicateurs de Lagrange, on maximise :
\[
  \mathcal{L}
    =
    \mathbf{a}_i^{H}
    \mathbf{R}_i
    \mathbf{a}_i
    -
    \lambda
    \left(
      \mathbf{a}_i^{H}\mathbf{a}_i
      -
      1
    \right).
\]

La condition d'optimalite donne :
\[
  \mathbf{R}_i\mathbf{a}_i
    =
    \lambda \mathbf{a}_i.
\]

Le maximum est donc atteint pour le vecteur propre associe a la plus grande
valeur propre de $\mathbf{R}_i$. C'est cette operation qui est calculee par la
decomposition en valeurs propres dans l'implementation.

\subsection{Mise a jour de la variance}

La densite de probabilite sur l'hypersphere complexe de dimension $M$ est
proportionnelle a :
\[
  p\left(
    \tilde{\mathbf{x}}_t
    \mid
    \mathbf{a}_i,\sigma_i^2
  \right)
  =
  \frac{1}
       {(\pi\sigma_i^2)^{M-1}}
  \exp\left(
    -
    \frac{
      1-\left|\mathbf{a}_i^{H}\tilde{\mathbf{x}}_t\right|^2
    }
    {\sigma_i^2}
  \right).
\]

En posant :
\[
  d_{i,t}^2
    =
    1
    -
    \left|
      \mathbf{a}_i^{H}\tilde{\mathbf{x}}_t
    \right|^2,
\]
le terme de $Q$ dependant de $\sigma_i^2$ s'ecrit :
\[
  Q(\sigma_i^2)
    =
    \sum_{t=1}^{T}
    \gamma_{i,t}
    \left[
      -(M-1)\ln(\sigma_i^2)
      -
      \frac{d_{i,t}^2}{\sigma_i^2}
    \right].
\]

On derive ensuite par rapport a $\sigma_i^2$ :
\[
  \frac{\partial Q}{\partial \sigma_i^2}
    =
    -
    \frac{
      (M-1)\sum_t \gamma_{i,t}
    }
    {\sigma_i^2}
    +
    \frac{
      \sum_t \gamma_{i,t}d_{i,t}^2
    }
    {(\sigma_i^2)^2}.
\]

En annulant cette derivee, on obtient :
\[
  (M-1)\sigma_i^2\sum_t\gamma_{i,t}
    =
    \sum_t\gamma_{i,t}d_{i,t}^2.
\]

D'ou :
\[
  \sigma_i^2
    =
    \frac{
      \sum_t \gamma_{i,t}d_{i,t}^2
    }
    {
      (M-1)\sum_t \gamma_{i,t}
    }.
\]

\end{document}
